\section{Introduction}
\label{sec:intro}

Networks trained on long sequences of tasks can lose the ability to learn new ones
\citep{dohare2024loss,lyle2023understanding}. Adam \citep{kingma2015adam} is a recurring suspect: its
second-moment estimate adapts slowly after an abrupt change, producing unusually large or unusually small
updates \citep{lyle2023understanding,dohare2024loss,ellis2024adam}. This has motivated interventions that
act only on the optimizer state, leaving the weights untouched: resetting both moments and the counter at
known boundaries \citep{asadi2023resetting}, resetting only the step counter \citep{ellis2024adam}, or
carrying and mixing second moments across tasks \citep{kang2024continual}.

An intervention on Adam's state changes the magnitude of the next updates as well as their direction.
If an intervention helps, the benefit may therefore come from ``repairing'' stale statistics or simply
from taking larger (or smaller) steps. \citet{ellis2024adam} already note that their counter reset
behaves like learning-rate annealing or warm-up. We ask a narrower diagnostic question:
\emph{in a setting where a loss of learning ability is actually observed on a comparable new task, does a
correctly bias-corrected optimizer-state intervention help beyond what matching its update scale explains?}

We answer it in a deliberately small, fully instrumented setting. All branches start from one checkpoint,
see the same batches and dropout masks, and are checked to differ from the untouched checkpoint only in
optimizer state at the moment of intervention. Each intervention is paired with an update-norm control
that uses its own optimizer state for the direction and the intervention's per-step, per-tensor norms
for the magnitude. Our findings are:

\begin{itemize}
\item \textbf{The phenomenon is present in our setting.} On two synthetic 50-task streams, the late
network's probe learning-curve area is lower than the early network's by \RTLateMinusEarlyAuc{}
(random teacher) and \LPLateMinusEarlyAuc{} (label permutation), with every seed negative
(Section~\ref{sec:phenomenon}).
\item \textbf{We found no support for a benefit of correctly corrected state-only interventions beyond
update scale.} A fresh optimizer state on the late weights did not recover the gap. The only intervention
with a consistent positive difference to \keep{}, resetting the first moment, exceeded its
update-norm control by less than our minimum interesting effect (Section~\ref{sec:interventions}).
This is absence of support at three seeds, not evidence of zero effect.
\item \textbf{The algebra of counter and moment policies explains two of the observations}
(Section~\ref{sec:analysis}). With $\eps=0$ a counter-only reset rescales every coordinate of a tensor
by the same factor, so it is a learning-rate schedule. With $\eps>0$ this equivalence is only
approximate. Exact bias correction of a reset second moment does not bound the update, because the kept
first moment can be large where the current gradient is small. These are elementary consequences of the
Adam update, stated for clarity rather than as new results.
\end{itemize}

We do not propose a new optimizer, do not claim to restore plasticity, and do not claim that features
stay fixed. The weights of every branch are updated after the intervention.
